\subsection{Preserving the full ordinary pair during the pro-three induction}

\begin{result}{E.151}{Proved}{The finite plus gate can retain ordinary compatibility at every coefficient prime}
We strengthen the conditional construction in E.148; the existence of
an infinite path through its finite plus equations remains undecided.
Put
\[
 T_n=\frac{(n+11)!}{3^{v_3((n+11)!)}},\qquad
 M_n=3^{n-1}T_n,\qquad B(m)=(m+1)^4\quad(n\ge2).
 \tag{E.151.1}
\]
Thus \(M_n\mid M_{n+1}\), \(v_3(M_n)=n-1\), and
\(M_n\) tends to zero in \(\widehat{\mathbb Z}\).
At level \(n-1\) retain a literal section in \(Q_{n-1}\),
unramified outside \(S\), and positive, pairwise coprime integers
\[
 a=5U^8,\qquad b=11V^8,\qquad c=16W^8,
 \qquad (UVW,2\cdot3\cdot5\cdot11)=1.
 \tag{E.151.2}
\]
The prime factors of \(UVW\) are distinct support primes congruent
to two modulo nine; the three supports are disjoint.
We require the following finite-stage invariants:
\begin{enumerate}
\item the abelianization of the section is
\((\kappa(a/c),\kappa(b/c))\) modulo \(3^{n-1}\);
\item at each support prime it agrees with the actual ordinary
point class having the prescribed integer valuation pattern, and at
three its abelianization agrees with the fixed point \(t_3=5/16\);
\item the curve \(aX^{M_n}+bY^{M_n}=cZ^{M_n}\) has a
unit-coordinate point over every finite \(\mathbb Q_p\);
\item with \(H_n(p)=2v_p(M_{n+1})+n+3\),
\[
 v_p(a+b-c)\ge H_n(p)\qquad(p\le B(M_n)).
 \tag{E.151.3}
\]
\end{enumerate}
Here \(S\) consists of three, infinity, and the support primes;
in particular two, five, and eleven belong to \(S\).
The rational-point section and \((a,b,c)=(5,11,16)\) initialize
these invariants at \(n=2\).

For these data the next stage exists if and only if the quotient
\(Q_n/B_n\) has a global lift over \(G_{\mathbb Q,S}\)
matching the old support point classes. Equivalently, for one of the
finite aligning-frame choices it satisfies the plus equations
(E.148.2). If it exists, we can obtain it by introducing between
nine and twenty-one new primes, all two modulo nine and arbitrarily
large, and by multiplying each of \(a,b,c\) by a product of their
\(M_n\)-th powers. The ordinary integer valuation at each newly
introduced prime is frozen thereafter.

If this construction has an infinite path, its limits give one
ordinary completed pair in \(\widehat{\mathbb Q^\times}^{,2}\)
of infinite support, together with one continuous full geometrically
pro-three section above its three-primary projection. Thus the
remaining condition for this construction is the compatible infinite
solution of the old-support plus equations; compatibility with the
other coefficient primes in the ordinary pair introduces no extra
finite-step condition. This is not a construction of an infinite
path, nor of sections in the other pro-prime components.
\end{result}
\begin{proof}
We give the simultaneous arithmetic and section construction.
All finite geometric groups, actual arithmetic actions and peripheral
frames are those of E.144--E.148.

\emph{The normalized weak lift.}
Suppose the quotient gate holds. By E.146 its \(B_n\)-obstruction
vanishes and it has a global weak \(Q_n\)-lift unramified outside
\(S\). Recall
\[
 B_n=\langle x^{3^{n-1}},y^{3^{n-1}}\rangle
       \simeq\mu_3^2.
 \tag{E.151.4}
\]
Twist by \(S\)-unit Kummer classes to match the old support point
classes and to make the two abelian valuation coordinates at three
zero. At a support prime \(p\), we have \(p\equiv-1\pmod3\),
including \(p=2\). The local \(H^1(\mu_3)\) is its valuation
line and the inclusion into \(H^1(\mu_{3^n})\) is injective,
since \(H^0(\mathbb Q_p,\mu_{3^{n-1}})=0\). Thus these
normalizations are compatible and give exactly the prescribed
integer abelian valuations modulo \(3^n\).
Over \(\mathbb Q\), a finite Kummer class unramified outside
\(S\) is determined by these valuations: the remaining rational
unit is \(\pm1\), which is a \(3^n\)-th power. Consequently
the normalized weak lift has abelianization precisely
\((a/c,b/c)\) modulo \(3^n\). We have not yet imposed its
ordinary unit digits at three.

\emph{Local product targets.}
Write \(m=M_n\), \(h=M_{n+1}\), and choose \(P\ge B(h)\)
larger than all old support primes. Let \(\mathcal T\) be the
finite set of primes at most \(P\). Choose unit solutions
\((X_p,Y_p,Z_p)\) of the exponent-\(m\) curve at these primes.
At \(p=2,5,11\), choose them all congruent to one modulo
\(4,5,11\), respectively. This follows directly from (E.151.3):
keep two coordinates equal to one and apply the sharpened Newton
estimate of E.102 to a variable with unit coefficient. Its derivative
has valuation \(v_p(m)\), whereas the error has valuation at least
\(2v_p(h)+n+3\). In particular the resulting root is one modulo
four at two and one modulo \(p\) at five and eleven.

At three choose a solution realizing the specific point \(t_3\):
\[
 X_3^m=t_3c/a,\qquad Y_3^m=(1-t_3)c/b,
 \qquad Z_3^m=1.
 \tag{E.151.5}
\]
Both right sides are \(m\)-th powers in \(\mathbb Q_3^\times\).
Indeed their three-primary Kummer classes vanish modulo
\(3^{n-1}\) by the preceding-stage hypothesis. They are units
congruent to one modulo three, by (E.151.2). On \(1+3\mathbb Z_3\)
the prime-to-three part \(T_n\) is an automorphism; the only
remaining condition is exactly divisibility by \(3^{n-1}\) in
the logarithm. Since \(m\) is even, we may change each root's
sign independently and take
\(X_3,Y_3,Z_3\equiv-1\pmod3\).

For each direction \(i\in\{a,b,c\}\), choose
\(k_i\in\{3,5,7\}\) with
\[
 2^{k_a}\equiv X_3,\quad 2^{k_b}\equiv Y_3,
 \quad 2^{k_c}\equiv Z_3\pmod9.
 \tag{E.151.6}
\]
This is possible because \(2^3,2^5,2^7\) are respectively
\(8,5,2\) modulo nine, the three unit residues negative modulo
three. Factor each local product target into \(k_i\) local units.
Away from three take the first \(k_i-1\) factors equal to one
and the last equal to the target. At three take the first factors
equal to two and the last equal to the target divided by
\(2^{k_i-1}\). Each individual factor is then two modulo nine.
At two, five and eleven, each factor is one modulo four, five,
and eleven, respectively. These choices involve no extraction
of a root of order \(k_i\) in a residue-field unit group.

Fix CRT precisions \(L_p\) large enough that replacing these
factors by congruent primes makes the new sum error at least
\(H_{n+1}(p)\) at every \(p\in\mathcal T\), and makes the
new two ratios at three agree with \(t_3,1-t_3\) modulo the
\(3^n\)-power subgroups. Explicitly it suffices to take
\[
 L_p\ge H_{n+1}(p)+n+3,
 \qquad L_p\ge v_p(h)+1,
 \tag{E.151.7}
\]
and increase \(L_3\) if necessary; polynomial evaluation on
integral units preserves these congruences. At three all coefficients
are units, and \(1+3^{n+1}\mathbb Z_3\) consists of
\(3^n\)-th powers, so the displayed lower bound already suffices.

\emph{The required field intersections.}
We choose the new primes successively. At a step, let
\((a_*,b_*,c_*)\) be the coefficients after the previously
chosen new primes have been inserted. They still have the forms
(E.151.2). For the next prime in direction \(a,b,c\), respectively,
put
\[
 w=c_*/b_*,\qquad w=c_*/a_*,\qquad w=-b_*/a_*.
 \tag{E.151.8}
\]
In addition to its CRT conditions and cusp Frobenius prescription,
we require \(w\) to be an \(h\)-th power modulo that prime.

We verify compatibility, instead of assuming that these conditions
are independent. The squareclasses of the three possible \(w\)
are \(11,5,-55\). For every cyclotomic field \(F\) containing
\(\mu_h\), these elements are squares, since \(220\mid h\).
No \(w\) is an odd-prime power in an abelian extension: its
valuation at five or eleven is \(1\) or \(-1\), so the real
pure odd-prime radical has that prime degree and is not normal.
If a root were in an abelian extension containing the necessary
roots of unity, its real root would be there too, a contradiction.
No \(w\) is a fourth power there either. Removing its rational
square factors, and \(\zeta_8\) in the negative case, would put
one of \(5^{1/4},11^{1/4},55^{1/4}\) in an abelian field.
These real quartic fields are nonnormal (Eisenstein at five or
eleven); the factor \(\sqrt5\) needed for the ratio in the
last case is already cyclotomic. This is the same explicit radical
argument as E.102, now with the private valuations at five and eleven.
It proves
\[
 [F(w^{1/h}):F]=h/2.
 \tag{E.151.9}
\]
Equivalently, for a chosen root \(z\), the polynomial
\(X^{h/2}-\sqrt w\) is irreducible over \(F\).
An automorphism \(\sigma\) of \(F/\mathbb Q\) extends fixing
some such \(z\) whenever it fixes \(\sqrt w\): use this
irreducible polynomial with its fixed coefficient. Our CRT prime
classes are one modulo \(220\), and so fix all three square roots.

Let \(L\) be the finite joint field of the quotient lift,
the actual action, and the peripheral frames, as in E.146.
It has a normal three-subgroup with quotient the scalar group
of order two. Enlarge the arithmetic frame field to include the
needed finite three-cyclotomic data. The independent Frattini inertia
vectors at five and eleven imply, by the finite Burnside basis
argument of E.146, that the geometric image over that frame field
is all of \(Q_n/B_n\). Hence above any prescribed arithmetic
Frobenius value there is an element whose quotient-section value
is the desired cusp value in the chosen direction.

Take \(F\) to be a cyclotomic field containing \(\mu_h\),
all CRT moduli so far, and \(L\cap\mathbb Q^{\rm ab}\).
The latter field is compatible with the old CRT prescriptions.
To see this explicitly, every abelian three-extension contained
in \(L\) is unramified at the old primes other than three:
their tame residue degrees \(p-1\) are prime to three, and
their wild inertia is pro-\(p\). Its three-part is therefore
contained in the three-power cyclotomic extension. Its only possible
quadratic subfield is \(\mathbb Q(\mu_3)\). Thus the prescribed
three-adic CRT class fixes exactly the needed abelian value.
Alternatively the residual scalar \(-1\) kills the whole geometric
Frattini quotient in the abelianization, so all abelian data factor
through the arithmetic frame field.

Choose the element on \(LF\) with that cyclotomic class and the
desired cusp quotient value. Over \(F\), \(LF/F\) is a
three-extension. Its intersection
\[
 J=LF\cap F(w^{1/h})
 \tag{E.151.10}
\]
is therefore cyclic of three-power degree. On the three-primary
radical-root torsor the prescribed automorphism has affine form
\(v\mapsto\chi v+b\), where \(\chi\equiv-1\pmod3\).
Since \(\chi-1\) is invertible, this affine transformation has
a unique fixed point at every finite three-power level.
More explicitly, start with a fixed-root extension of \(\sigma\)
to \(F(w^{1/h})\). Conjugating it by the three-Sylow subgroup
of its cyclic Kummer group changes its restriction to \(J\)
by multiplication by \(1-\chi\), a bijection. It can therefore
be made to agree with the prescribed element on \(LF\), while
still fixing the conjugate radical root. The two restrictions
define an actual automorphism of the compositum. Chebotarev
\citep[Theorem~8.31, p.~147]{MilneANT2020} supplies arbitrarily
large primes with this Frobenius conjugacy class. A fixed radical
root yields an \(h\)-th power in the residue field, and cusp
matching is retained up to the geometric conjugacy used throughout.

In choosing each later prime, impose in addition that it is one
modulo every new prime chosen earlier in this stage. These extra
cyclotomic conditions are independent CRT conditions. They do not
alter the compatibility argument: the earlier new primes are
\(-1\) modulo three, the quotient/frame field is unramified there,
and the radical argument (E.151.9)--(E.151.10) uses the current
coefficients. This completes the sequential prime selection.

\emph{The next arithmetic invariants.}
Write the products of selected primes in the three directions as
\(P_a,P_b,P_c\), and set
\[
 a'=aP_a^m,
 \qquad b'=bP_b^m,
 \qquad c'=cP_c^m.
 \tag{E.151.11}
\]
Their forms and disjoint supports persist because \(8\mid m\).
At all old primes at most \(P\), the product targets solve the
old curve exactly. The CRT precisions give the new sum-error
invariant. Newton at \((1,1,1)\), in a variable with unit
coefficient, gives a unit solution for exponent \(h\), since
the error valuation is greater than \(2v_p(h)\).
At a newly introduced prime in direction \(a\), its introduction
makes the \(a\)-coefficient zero modulo that prime; (E.151.8)
provides the needed \(h\)-th power \(c_*/b_*\).
All later new primes are one modulo it, so this residue condition
survives to the end of the stage. Set the other coordinate equal
to one and use simple-root Hensel to obtain a unit-coordinate
solution of the exponent-\(h\) curve. The directions \(b,c\)
use \(c_*/a_*\) and \(-b_*/a_*\) in exactly the displayed
equation, respectively. The derivative is a unit because every
new prime exceeds \(B(h)\).
At every remaining prime the coefficients are units and the prime
exceeds \(B(h)\); the explicit diagonal-curve Weil bound proved
in E.31 gives a unit solution. All arithmetic invariants follow.

\emph{The next section and necessity of the gate.}
Twist the normalized weak lift by the Kummer classes of the new
primes in the corresponding peripheral directions, with coefficient
\(m/3^{n-1}=T_n\) in \(\mathbb F_3\). Its abelianization is
exactly (E.151.11) modulo \(3^n\). Old support matching is
unchanged, since every new prime is a local unit and those unit
Kummer groups vanish at the old primes. At three (E.151.5)--(E.151.7)
give the prescribed abelian point class. At a new prime, the
quotient Frobenius is its prescribed cusp value and the inertia
is the appropriate peripheral \(m\)-th power. The vanishing
\(H^1_{\rm nr}(\mathbb Q_q,B_n)=0\), because \(q-1\) is a
three-adic unit, gives the actual ordinary local point class by
E.146. Any ordinary point with the same integer valuation pattern
has the same full pro-three local class by A.31; hence later
arithmetic refinements of its units preserve this fixed target.
This verifies every finite-stage section invariant.

Conversely, an extension with ordinary images at its new support
primes must pass the quotient gate by E.144. In particular the
construction does not remove that necessary condition by adding
auxiliary ramification. The equivalence with the finite plus
equations is E.148, with all aligning frames retained.

\emph{The conditional limit.}
Suppose that an infinite sequence of these choices exists.
The literal finite sections give a continuous section into
\(F=\varprojlim Q_n\), as in E.148. The three products with
exponents \(M_n\) converge in the full rational power completion,
because every fixed positive integer divides all sufficiently late
\(M_n\). At a fixed prime, coefficient valuations eventually
stabilize; the later ratios are \(M_n\)-th powers of local units.
These tend uniformly to one in every finite quotient of the compact
unit group, so the three coefficients converge to nonzero ordinary
local values. The error invariant eventually applies at that prime
with arbitrarily large precision, and their limits satisfy
\(a_p+b_p=c_p\). Openness of local power subgroups identifies
their ratios with the localizations of the full completed products,
exactly as in E.102. Thus \(t_p=a_p/c_p\) is an ordinary local
point at every finite prime. The positive real signs are represented
by \(t_\infty=1/2\).
Every stage introduces new support in each direction with nonzero
integer valuation \(M_n\), so the completed pair has infinite
support. Its three-primary abelianization is that of the single
limiting section by coordinatewise compatibility. This proves the
conditional assertion, while leaving the infinite plus-gate path
and all other pro-prime sections undecided.
\end{proof}
